\chapter{Legal, Social, Ethical and Professional Issues}\label{ch:lsepi}

This project builds security-critical infrastructure for a widely used messaging application, and
it does so by modifying deployed open-source software; both facts carry legal, social, ethical, and
professional obligations. This chapter sets out those obligations and how they were discharged
during the work. It considers the legal position of the code and the data
(\cref{sec:lsepi-legal}), the social and ethical implications of strengthening, and of being
seen to strengthen, end-to-end-encrypted messaging (\cref{sec:lsepi-social}), and the
professional standards, specifically the BCS Code of Conduct~\cite{bcs-code} and the IET Rules of
Conduct~\cite{iet-conduct}, against which the work was conducted (\cref{sec:lsepi-prof}).

% ============================================================================
\section{Legal Issues}\label{sec:lsepi-legal}

Intellectual property and licensing are the most immediate legal consideration. The prototype
extends two Signal codebases, libsignal and Signal Desktop, both released under the GNU
Affero General Public License v3 (AGPL-3.0)~\cite{libsignal,signal-desktop}. The project is
distributed on the basis that modifications to the AGPL-covered Signal code remain under AGPL-3.0:
their source must remain available, and, distinctively for the AGPL, that obligation extends to
users who interact with a modified version over a network, not only to those who receive a binary.
Whether a particular separately developed component is itself a derivative work can depend on the
facts, the jurisdiction, and the mode of combination and distribution. The prototype is a
research artefact rather than a deployed service, so no public network offering has been made; but
the licence terms have been respected throughout: the modifications are kept in a fork whose
source is available, the upstream copyright and licence notices are preserved, and any future
deployment of the modified Signal code would carry the same AGPL-3.0 obligations, including the
network-interaction source requirement. The independent Rust \texttt{duet} crate depends on
permissively licensed libraries, including \texttt{sha2} and \texttt{hkdf} (MIT/Apache-2.0)
and \texttt{curve25519-dalek} and \texttt{ed25519-dalek} (BSD-3-Clause), all compatible with that
outcome. All third-party work (specifications, prior systems, and code) is attributed in the
references.

A second consideration is export control. The work implements cryptographic software, which is
regulated under the Wassenaar Arrangement and national export regimes; libsignal itself is
published under the corresponding open-source exception for publicly available cryptographic source
code. The contribution composes standardized primitives
(the RFC~9381 ECVRF, RFC~5869 HKDF, and Ed25519) without introducing new cryptographic algorithms, and is released as open, published source. This
thesis does not determine the software's export-control classification: public availability may be
relevant, but the applicable control lists and end-use rules require a case-specific assessment.

Data protection is the third. Key transparency is, by construction, a data-minimizing design: the
directory stores public keys and pseudonymous labels, never message content, and the VRF-derived
global labels (\cref{ssec:design-labels}) deliberately prevent an observer from enumerating which
users are present. The benchmark corpus used generated identifiers not linked to identifiable
individuals. Subject to that assumption, the benchmark data were not personal data and no
obligations under the UK GDPR were engaged by the measurements, while the design itself advances the
privacy interests that data-protection regimes exist to protect. Screenshots and operational logs that contain
identifiable names or paths would require separate review before release. A
production deployment would inherit one governance duty the prototype does not: any operational
logging of warnings or audit events must itself be data-minimized and access-controlled, so that
the alerting layer does not re-expose the label-to-user metadata the pseudorandom labels are
designed to hide.

% ============================================================================
\section{Social and Ethical Issues}\label{sec:lsepi-social}

The social value of the work is a direct contribution to public wellbeing and security:
authenticated key exchange is what makes end-to-end encryption meaningful, and the people who most
depend on it (journalists, activists, and others at risk) are exactly those whom an
undetected key substitution would endanger. Strengthening the automatic detection of such attacks,
and giving both conversation partners an in-band warning, serves the public interest that the BCS
code places first.

That same importance imposes a duty of honesty about what the work does and does not achieve,
because a security tool trusted beyond its guarantees is worse than none. The thesis is explicit
about its limitations: the prototype carries two proof-of-concept stand-ins for live protocol state
(\cref{sec:impl-strategy}), it is a single-machine prototype rather than a distributed deployment,
and the boundary case of key extraction is out of scope (\cref{ch:threat}). The user-facing design
embodies the same caution: an unreachable server raises a calm, amber notice rather than the red
alarm reserved for a detected attack (\cref{ssec:impl-verifier}), precisely so that the system does
not ``cry wolf'' and erode the trust its warnings depend on. The design is open and
specification-grounded rather than security-through-obscurity, and its correctness is evidenced by
reproducible tests against published vectors (\cref{sec:impl-tests}) rather than asserted.

A dual-use dimension is unavoidable: strong messaging security protects everyone, including those
who would misuse it. This project treats defensive integrity infrastructure of this kind as serving the public good,
since weakening it to disadvantage bad actors would also weaken protection for the far larger
population of ordinary and at-risk users; the professional codes cited in \cref{sec:lsepi-prof}
support this position. Finally, the work followed responsible-disclosure norms in the small: the one
security-relevant weakness found during integration, a per-lookup head-binding gap, was in
the author's own prototype and was fixed (\cref{sec:impl-challenges}); no vulnerability in deployed
Signal software was discovered, but had one been, the appropriate course would have been private
disclosure to Signal before any publication.

A warning mechanism also has direct effects on the people who see it, which a responsible design must
weigh. Because the channel can be written by a peer, a malicious or compromised contact can raise a
spurious warning; at scale, frequent or unexplained alerts risk \emph{warning fatigue}, in which users
learn to dismiss the very signal meant to protect them (\cref{sec:disc-limitations}). The design
mitigates this by distinguishing a genuine integrity alert (red) from mere unavailability (amber) and
by deriving the warning from cryptographic evidence rather than heuristics, but the wording, frequency,
and explicability of the alert are a human-factors problem this work does not claim to have solved, and
getting them wrong could erode trust rather than build it. The stakes are highest for at-risk users:
in an intimate-partner-abuse setting a key change, or a warning about one, can carry safety
implications beyond the cryptographic question, and any deployment should treat the accessibility and
clear explanation of these alerts, and a safe path to seek help, as first-order requirements. On the
conduct of the research itself, finally, all adversarial functionality was exercised only in an
isolated local test harness against synthetic accounts, with no live Signal interception and no real
users (\cref{ssec:impl-demo}).

% ============================================================================
\section{Professional Issues and Codes of Conduct}\label{sec:lsepi-prof}

The work was carried out in line with the BCS Code of Conduct~\cite{bcs-code} and the IET Rules of
Conduct~\cite{iet-conduct}; the following sets out concretely how.

The BCS code is organized under four principles. \emph{Public interest} (due regard for
security, privacy, and wellbeing) is the project's subject matter, addressed above.
\emph{Professional competence and integrity} shaped the technical choices: established, audited
cryptographic libraries were used in preference to bespoke implementations, and the one primitive
implemented from scratch, the ECVRF, was validated against the RFC~9381 test vectors rather than
trusted on inspection (\cref{ssec:impl-core-crypto}); claims are held to what the evidence supports,
and an unfamiliar literature was studied in order to work competently within it. \emph{Duty to the
relevant authority} was met through honest, regular engagement with the project supervisor and
adherence to the university's academic-integrity requirements, including full attribution of
sources and a clear demarcation of the proof-of-concept stand-ins from verified results.
\emph{Duty to the profession} is served by not overstating the contribution and by releasing the
work openly for scrutiny.

The IET Rules of Conduct, which adopt the Engineering Council and Royal Academy of Engineering
Statement of Ethical Principles (Rule~2), point the same way. Their emphasis on \emph{accuracy and
rigour} is reflected in the exhaustive bounded and invariant-oriented testing and the reproducibility measures
(\cref{sec:impl-tests,sec:impl-repro}); their requirement to work within one's competence and to
keep knowledge current (Rules~5--6) is reflected in the grounding of every construction in a named
specification, and their insistence that a professional opinion state its limitations clearly
(Rule~26) is reflected in the explicit scoping of the claims throughout. Above all, both codes
place the public good first, which for a project whose entire purpose is to make a widely used
private-messaging system harder to attack is not an external constraint but the point of the
exercise.